\chapter{Marco teórico}\label{chap:marco}

Este capítulo establece los fundamentos conceptuales necesarios para el proyecto \textit{BasketballAnalysis}. 
Primero se describen conceptos clave del baloncesto cuantitativo (posesiones, eficiencia, cuatro factores, 
alineaciones y contexto competitivo). Después, se presenta el armazón de aprendizaje automático empleado, 
incluyendo formulación del problema, validación temporal, modelos elegidos y consideraciones de calibración e interpretabilidad.

\section{Baloncesto: fundamentos y métricas}\label{sec:baloncesto}

\subsection*{Posesiones y ritmo}
El análisis comparativo exige normalizar por volumen de posesiones. Una aproximación habitual al conteo de posesiones de un equipo es:
\begin{equation}
\mathrm{Poss} \approx \mathrm{FGA} + 0.44\,\mathrm{FTA} - \mathrm{OREB} + \mathrm{TOV}.
\end{equation}
El \textbf{ritmo} o \textbf{PACE} expresa posesiones por 48 minutos y permite comparar equipos que juegan a velocidades distintas.

\subsection*{Eficiencia de tiro}
La \textbf{eFG\%} ajusta el valor del triple:
\begin{equation}
\mathrm{eFG\%} = \frac{\mathrm{FGM} + 0.5 \cdot \mathrm{FG3M}}{\mathrm{FGA}}.
\end{equation}
La \textbf{TS\%} aproxima la eficiencia de anotación global incorporando tiros libres:
\begin{equation}
\mathrm{TS\%} = \frac{\mathrm{PTS}}{2\,(\mathrm{FGA} + 0.44\,\mathrm{FTA})}.
\end{equation}

\subsection*{Ratings por 100 posesiones}
\begin{equation}
\mathrm{ORtg} = 100 \cdot \frac{\mathrm{PTS}}{\mathrm{Poss}}, \qquad
\mathrm{DRtg} = 100 \cdot \frac{\mathrm{PTS\ concedidos}}{\mathrm{Poss}}, \qquad
\mathrm{NetRtg} = \mathrm{ORtg} - \mathrm{DRtg}.
\end{equation}
Estos indicadores permiten separar contribuciones ofensivas y defensivas normalizadas.

\subsection*{Cuatro factores}
Según la descomposición clásica, el rendimiento ofensivo se resume en cuatro pilares:
(i) eficacia de tiro (eFG\%), (ii) pérdidas (TOV\%), (iii) rebote ofensivo (OREB\%) y (iv) tiros libres (FTr).
Analizar su contribución relativa ayuda a explicar variaciones en \textit{NetRtg} entre equipos, alineaciones y contextos.

\subsection*{Alineaciones, on/off y estabilidad}
La composición de quintetos modula los espacios de tiro, la creación y la contención defensiva.
Las métricas \textit{on/off} estiman el impacto de un jugador comparando el rendimiento del equipo cuando está en pista
frente a cuando no lo está. No obstante, requieren \textit{shrinkage} o regularización por tamaño muestral para evitar conclusiones espurias.
La \textbf{estabilidad de rotación} (varianza de quintetos utilizados) es una pista útil de consistencia táctica.

\subsection*{Contexto competitivo}
El lugar del partido (local/visitante), los días de descanso, los \textit{back-to-back} y el calendario reciente condicionan el rendimiento.
Las segmentaciones temporales (meses, pre/post All-Star) capturan cambios de tendencia y ajustes de plantilla.

\section{Modelos y aprendizaje automático}\label{sec:modelos}

\subsection*{Formulación del problema}
El proyecto incluye estimadores de distinta naturaleza. El \textbf{Estimador 1} se plantea como clasificación binaria (W/L por partido).
Otros estimadores pueden formularse como regresión (diferencial de puntos) o agregación/proyección (victorias de temporada).

\subsection*{Representación de datos y prevención de \textit{leakage}}
Se construye una tabla de \textbf{partido–equipo}, donde cada fila representa un equipo en un partido dado.
Todas las variables (\textit{features}) se calculan usando únicamente información disponible \emph{antes} del salto inicial.
Se emplean diferencias simétricas frente al rival (\texttt{Feature\_Diff = Equipo - Rival}) para centrar el modelo en ventajas relativas.

\subsection*{Validación temporal}
La evaluación se realiza con \textbf{backtesting \textit{walk-forward}}: se entrena en el histórico hasta un punto temporal $T$ y se valida en el intervalo $(T, T+\Delta]$.
Este esquema respeta la causalidad y evita fugas de información propias de particiones aleatorias.

\subsection*{Modelos utilizados}
\paragraph{Random Forest (RF).}
Bosque de árboles con \textit{bagging}. Ventajas: robustez, poco \textit{tuning} y buena línea base.
Parámetros prácticos: número de árboles alto, profundidad moderada, tamaño mínimo de hoja suficiente para evitar sobreajuste.

\paragraph{XGBoost (XGB).}
\textit{Gradient boosting} con regularización explícita.
Capta no linealidades e interacciones de forma eficiente y suele optimizar métricas como LogLoss con buen sesgo-varianza.
Parámetros clave: \texttt{learning\_rate}, \texttt{max\_depth}, \texttt{subsample}, \texttt{colsample\_bytree}, \texttt{min\_child\_weight}, y regularización L1/L2.

\subsection*{Ponderación temporal y calibración}
Las observaciones más recientes se ponderan más (decaimiento exponencial por días).
La salida probabilística se \textbf{calibra} (Platt o isotónica) para que los porcentajes predichos coincidan con frecuencias observadas,
lo cual es crucial en toma de decisiones.

\subsection*{Interpretabilidad}
Se emplean importancias por permutación y, en su caso, \textit{SHAP values} para explicar contribuciones globales y locales.
La interpretabilidad permite auditar el comportamiento del modelo y detectar variables espurias o inestables.

\subsection*{Despliegue y control de calidad}
En producción, el flujo actualiza las \textit{features} hasta la víspera, infiere para los partidos del día y genera reportes de rendimiento,
calibración y deriva (cambios de distribución) para mantenimiento preventivo.
